% =====================================================================
%  Anatomía de un LLM — Fase VI — Módulo 24
%  Formular y resolver problemas de optimización
%  Compilar: latexmk -r ../latexmkrc M24-formulacion.tex
% =====================================================================
\documentclass[11pt,a4paper,twoside,openany]{book}
\usepackage{../estilo-apuntes}
\usepackage{makecell}

\newcommand{\ApuntesCurso}{Anatomía de un LLM}
\newcommand{\ApuntesModulo}{Módulo 24}

% Notación propia del módulo
\DeclareMathOperator*{\argmin}{arg\,min}
\DeclareMathOperator*{\argmax}{arg\,max}
\newcommand{\Z}{\mathbb{Z}}
\newcommand{\st}{\text{s.a.}\ }
\newcommand{\T}{^{\mathsf{T}}}
\newcommand{\Lag}{\mathcal{L}}
\newcommand{\clip}{\operatorname{clip}}
\newcommand{\Prob}{\mathbb{P}}
\newcommand{\eur}{\text{\,€}}
\makeatletter\g@addto@macro\UrlBreaks{\do\-\do\/}\makeatother

\hypersetup{pdftitle={Módulo 24. Formular y resolver problemas de optimización (resumen)},pdfauthor={Anatomía de un LLM}}

\makeatletter
\newlabel{ej:piensos}{{1.7}{2}{Caso de estudio: la fábrica de piensos}{ejemplo.7}{}}
\newlabel{ej:piensos@cref}{{[ejemplo][7][]1.7}{[1][2][]2}}
\newlabel{tab:piensos}{{1.1}{3}{Datos del caso de la fábrica de piensos}{table.1.1}{}}
\newlabel{tab:piensos@cref}{{[table][1][1]1.1}{[1][2][]3}}
\newlabel{fig:piensos}{{1.1}{3}{Región factible del caso de los piensos (sombreada) y recta de nivel óptima}{figure.1.1}{}}
\newlabel{fig:piensos@cref}{{[figure][1][1]1.1}{[1][3][]3}}
\newlabel{fig:anatomia}{{1.2}{4}{Arquitectura de un optimizador con IA. Las flechas continuas forman el ciclo de decisión; las discontinuas, el de aprendizaje y supervisión}{figure.1.2}{}}
\newlabel{fig:anatomia@cref}{{[figure][2][1]1.2}{[1][3][]4}}
\newlabel{ej:casos-reales}{{1.9}{4}{Dos sistemas publicados}{ejemplo.9}{}}
\newlabel{ej:casos-reales@cref}{{[ejemplo][9][]1.9}{[1][4][]4}}
\newlabel{obs:errores1}{{1.11}{5}{Errores frecuentes al modelizar}{observacion.11}{}}
\newlabel{obs:errores1@cref}{{[observacion][11][]1.11}{[1][5][]5}}
\newlabel{fig:arbol}{{1.3}{6}{Árbol de decisión simplificado para elegir la familia de métodos. Los problemas con parámetros inciertos o decisiones secuenciales se tratan en el \cref {tema:6} y en los módulos posteriores}{figure.1.3}{}}
\newlabel{fig:arbol@cref}{{[figure][3][1]1.3}{[1][5][]6}}
\newlabel{eje:1-elementos}{{1.1}{6}{}{ejercicio.1.1}{}}
\newlabel{eje:1-elementos@cref}{{[ejercicio][1][1]1.1}{[1][6][]6}}
\newlabel{eje:1-mapa}{{1.2}{7}{}{ejercicio.1.2}{}}
\newlabel{eje:1-mapa@cref}{{[ejercicio][2][1]1.2}{[1][6][]7}}
\newlabel{eje:1-abs}{{1.3}{7}{}{ejercicio.1.3}{}}
\newlabel{eje:1-abs@cref}{{[ejercicio][3][1]1.3}{[1][7][]7}}
\newlabel{eje:1-correos}{{1.4}{7}{}{ejercicio.1.4}{}}
\newlabel{eje:1-correos@cref}{{[ejercicio][4][1]1.4}{[1][7][]7}}
\newlabel{eje:1-piensos}{{1.5}{7}{}{ejercicio.1.5}{}}
\newlabel{eje:1-piensos@cref}{{[ejercicio][5][1]1.5}{[1][7][]7}}
\newlabel{eje:1-minimax}{{1.6}{7}{}{ejercicio.1.6}{}}
\newlabel{eje:1-minimax@cref}{{[ejercicio][6][1]1.6}{[1][7][]7}}
\newlabel{eje:1-desviaciones}{{1.7}{8}{}{ejercicio.1.7}{}}
\newlabel{eje:1-desviaciones@cref}{{[ejercicio][7][1]1.7}{[1][7][]8}}
\newlabel{eje:1-diseno}{{1.8}{8}{}{ejercicio.1.8}{}}
\newlabel{eje:1-diseno@cref}{{[ejercicio][8][1]1.8}{[1][8][]8}}
\newlabel{ej:sombra-piensos}{{2.8}{10}{Precios sombra en la fábrica de piensos}{ejemplo.8}{}}
\newlabel{ej:sombra-piensos@cref}{{[ejemplo][8][]2.8}{[1][10][]10}}
\newlabel{ej:cervecera}{{2.12}{12}{Caso de estudio: plan semanal de una cervecera}{ejemplo.12}{}}
\newlabel{ej:cervecera@cref}{{[ejemplo][12][]2.12}{[1][12][]12}}
\newlabel{obs:errores2}{{2.13}{13}{Errores frecuentes en PL y PLEM}{observacion.13}{}}
\newlabel{obs:errores2@cref}{{[observacion][13][]2.13}{[1][13][]13}}
\newlabel{eje:2-holgura}{{2.1}{14}{}{ejercicio.2.1}{}}
\newlabel{eje:2-holgura@cref}{{[ejercicio][1][2]2.1}{[1][14][]14}}
\newlabel{eje:2-dual}{{2.2}{14}{}{ejercicio.2.2}{}}
\newlabel{eje:2-dual@cref}{{[ejercicio][2][2]2.2}{[1][14][]14}}
\newlabel{eje:2-outdoors}{{2.3}{14}{}{ejercicio.2.3}{}}
\newlabel{eje:2-outdoors@cref}{{[ejercicio][3][2]2.3}{[1][14][]14}}
\newlabel{eje:2-bb}{{2.4}{14}{}{ejercicio.2.4}{}}
\newlabel{eje:2-bb@cref}{{[ejercicio][4][2]2.4}{[1][14][]14}}
\newlabel{eje:2-mochila}{{2.5}{15}{}{ejercicio.2.5}{}}
\newlabel{eje:2-mochila@cref}{{[ejercicio][5][2]2.5}{[1][14][]15}}
\newlabel{eje:2-logica}{{2.6}{15}{}{ejercicio.2.6}{}}
\newlabel{eje:2-logica@cref}{{[ejercicio][6][2]2.6}{[1][15][]15}}
\newlabel{eje:2-bigm}{{2.7}{15}{}{ejercicio.2.7}{}}
\newlabel{eje:2-bigm@cref}{{[ejercicio][7][2]2.7}{[1][15][]15}}
\newlabel{eje:2-concava}{{2.8}{15}{}{ejercicio.2.8}{}}
\newlabel{eje:2-concava@cref}{{[ejercicio][8][2]2.8}{[1][15][]15}}
\newlabel{ej:vrptw}{{3.5}{17}{Caso de estudio: reparto de barriles con ventanas horarias}{ejemplo.5}{}}
\newlabel{ej:vrptw@cref}{{[ejemplo][5][]3.5}{[1][17][]17}}
\newlabel{ej:jobshop}{{3.9}{19}{Caso de estudio: taller mecánico}{ejemplo.9}{}}
\newlabel{ej:jobshop@cref}{{[ejemplo][9][]3.9}{[1][19][]19}}
\newlabel{obs:errores3}{{3.10}{19}{Errores frecuentes en problemas combinatorios}{observacion.10}{}}
\newlabel{obs:errores3@cref}{{[observacion][10][]3.10}{[1][19][]19}}
\newlabel{eje:3-contar}{{3.1}{20}{}{ejercicio.3.1}{}}
\newlabel{eje:3-contar@cref}{{[ejercicio][1][3]3.1}{[1][20][]20}}
\newlabel{eje:3-ruta}{{3.2}{20}{}{ejercicio.3.2}{}}
\newlabel{eje:3-ruta@cref}{{[ejercicio][2][3]3.2}{[1][20][]20}}
\newlabel{eje:3-2opt}{{3.3}{20}{}{ejercicio.3.3}{}}
\newlabel{eje:3-2opt@cref}{{[ejercicio][3][3]3.3}{[1][20][]20}}
\newlabel{eje:3-cuñas}{{3.4}{20}{}{ejercicio.3.4}{}}
\newlabel{eje:3-cuñas@cref}{{[ejercicio][4][3]3.4}{[1][20][]20}}
\newlabel{eje:3-johnson}{{3.5}{21}{}{ejercicio.3.5}{}}
\newlabel{eje:3-johnson@cref}{{[ejercicio][5][3]3.5}{[1][21][]21}}
\newlabel{eje:3-bigm}{{3.6}{21}{}{ejercicio.3.6}{}}
\newlabel{eje:3-bigm@cref}{{[ejercicio][6][3]3.6}{[1][21][]21}}
\newlabel{eje:3-cota}{{3.7}{21}{}{ejercicio.3.7}{}}
\newlabel{eje:3-cota@cref}{{[ejercicio][7][3]3.7}{[1][21][]21}}
\newlabel{ej:despacho}{{4.7}{24}{Despacho económico de tres grupos térmicos}{ejemplo.7}{}}
\newlabel{ej:despacho@cref}{{[ejemplo][7][]4.7}{[1][23][]24}}
\newlabel{obs:errores4}{{4.9}{25}{Errores frecuentes en optimización continua}{observacion.9}{}}
\newlabel{obs:errores4@cref}{{[observacion][9][]4.9}{[1][24][]25}}
\newlabel{eje:4-dcp}{{4.1}{25}{}{ejercicio.4.1}{}}
\newlabel{eje:4-dcp@cref}{{[ejercicio][1][4]4.1}{[1][25][]25}}
\newlabel{eje:4-kkt}{{4.2}{25}{}{ejercicio.4.2}{}}
\newlabel{eje:4-kkt@cref}{{[ejercicio][2][4]4.2}{[1][25][]25}}
\newlabel{eje:4-despacho2}{{4.3}{26}{}{ejercicio.4.3}{}}
\newlabel{eje:4-despacho2@cref}{{[ejercicio][3][4]4.3}{[1][26][]26}}
\newlabel{eje:4-agua}{{4.4}{26}{}{ejercicio.4.4}{}}
\newlabel{eje:4-agua@cref}{{[ejercicio][4][4]4.4}{[1][26][]26}}
\newlabel{eje:4-sensib}{{4.5}{26}{}{ejercicio.4.5}{}}
\newlabel{eje:4-sensib@cref}{{[ejercicio][5][4]4.5}{[1][26][]26}}
\newlabel{eje:4-dual}{{4.6}{26}{}{ejercicio.4.6}{}}
\newlabel{eje:4-dual@cref}{{[ejercicio][6][4]4.6}{[1][26][]26}}
\newlabel{eje:4-primer-orden}{{4.7}{26}{}{ejercicio.4.7}{}}
\newlabel{eje:4-primer-orden@cref}{{[ejercicio][7][4]4.7}{[1][26][]26}}
\newlabel{ej:planta}{{5.6}{29}{Caso de estudio: distribución en planta}{ejemplo.6}{}}
\newlabel{ej:planta@cref}{{[ejemplo][6][]5.6}{[1][29][]29}}
\newlabel{fig:planta}{{5.1}{30}{Dos distribuciones Pareto-óptimas del caso de la planta. A la izquierda, mínimo recorrido; a la derecha, soldadura y pintura separadas por expedición}{figure.5.1}{}}
\newlabel{fig:planta@cref}{{[figure][1][5]5.1}{[1][30][]30}}
\newlabel{obs:errores5}{{5.7}{30}{Errores frecuentes con metaheurísticas}{observacion.7}{}}
\newlabel{obs:errores5@cref}{{[observacion][7][]5.7}{[1][30][]30}}
\newlabel{eje:5-aceptacion}{{5.1}{31}{}{ejercicio.5.1}{}}
\newlabel{eje:5-aceptacion@cref}{{[ejercicio][1][5]5.1}{[1][31][]31}}
\newlabel{eje:5-frentes}{{5.2}{31}{}{ejercicio.5.2}{}}
\newlabel{eje:5-frentes@cref}{{[ejercicio][2][5]5.2}{[1][31][]31}}
\newlabel{eje:5-crowding}{{5.3}{31}{}{ejercicio.5.3}{}}
\newlabel{eje:5-crowding@cref}{{[ejercicio][3][5]5.3}{[1][31][]31}}
\newlabel{eje:5-apoyado}{{5.4}{31}{}{ejercicio.5.4}{}}
\newlabel{eje:5-apoyado@cref}{{[ejercicio][4][5]5.4}{[1][31][]31}}
\newlabel{eje:5-ox}{{5.5}{31}{}{ejercicio.5.5}{}}
\newlabel{eje:5-ox@cref}{{[ejercicio][5][5]5.5}{[1][31][]31}}
\newlabel{eje:5-pesos}{{5.6}{31}{}{ejercicio.5.6}{}}
\newlabel{eje:5-pesos@cref}{{[ejercicio][6][5]5.6}{[1][31][]31}}
\newlabel{eje:5-eps}{{5.7}{31}{}{ejercicio.5.7}{}}
\newlabel{eje:5-eps@cref}{{[ejercicio][7][5]5.7}{[1][31][]31}}
\newlabel{eje:5-metropolis}{{5.8}{32}{}{ejercicio.5.8}{}}
\newlabel{eje:5-metropolis@cref}{{[ejercicio][8][5]5.8}{[1][32][]32}}
\newlabel{ej:fresca}{{6.10}{35}{Cerveza sin pasteurizar con demanda predicha}{ejemplo.10}{}}
\newlabel{ej:fresca@cref}{{[ejemplo][10][]6.10}{[1][35][]35}}
\newlabel{obs:errores6}{{6.11}{36}{Errores frecuentes al combinar predicción y decisión}{observacion.11}{}}
\newlabel{obs:errores6@cref}{{[observacion][11][]6.11}{[1][36][]36}}
\newlabel{eje:6-fugas}{{6.1}{37}{}{ejercicio.6.1}{}}
\newlabel{eje:6-fugas@cref}{{[ejercicio][1][6]6.1}{[1][37][]37}}
\newlabel{eje:6-newsvendor}{{6.2}{37}{}{ejercicio.6.2}{}}
\newlabel{eje:6-newsvendor@cref}{{[ejercicio][2][6]6.2}{[1][37][]37}}
\newlabel{eje:6-spo}{{6.3}{37}{}{ejercicio.6.3}{}}
\newlabel{eje:6-spo@cref}{{[ejercicio][3][6]6.3}{[1][37][]37}}
\newlabel{eje:6-vss}{{6.4}{37}{}{ejercicio.6.4}{}}
\newlabel{eje:6-vss@cref}{{[ejercicio][4][6]6.4}{[1][37][]37}}
\newlabel{eje:6-robusta}{{6.5}{37}{}{ejercicio.6.5}{}}
\newlabel{eje:6-robusta@cref}{{[ejercicio][5][6]6.5}{[1][37][]37}}
\newlabel{eje:6-saa}{{6.6}{38}{}{ejercicio.6.6}{}}
\newlabel{eje:6-saa@cref}{{[ejercicio][6][6]6.6}{[1][37][]38}}
\newlabel{eje:6-fractil}{{6.7}{38}{}{ejercicio.6.7}{}}
\newlabel{eje:6-fractil@cref}{{[ejercicio][7][6]6.7}{[1][38][]38}}
\newlabel{eje:6-maldicion}{{6.8}{38}{}{ejercicio.6.8}{}}
\newlabel{eje:6-maldicion@cref}{{[ejercicio][8][6]6.8}{[1][38][]38}}
\newlabel{ap:soluciones}{{A}{39}{Soluciones de los ejercicios}{appendix.A}{}}
\newlabel{ap:soluciones@cref}{{[appendix][1][2147483647]A}{[1][39][]39}}
\newlabel{sol:eje:1-elementos}{{A.1}{39}{Tema 1}{section*.3}{}}
\newlabel{sol:eje:1-elementos@cref}{{[subappendix][1][2147483647,1]A.1}{[1][39][]39}}
\newlabel{sol:eje:1-mapa}{{A.1}{39}{Tema 1}{section*.4}{}}
\newlabel{sol:eje:1-mapa@cref}{{[subappendix][1][2147483647,1]A.1}{[1][39][]39}}
\newlabel{sol:eje:1-abs}{{A.1}{40}{Tema 1}{section*.5}{}}
\newlabel{sol:eje:1-abs@cref}{{[subappendix][1][2147483647,1]A.1}{[1][40][]40}}
\newlabel{sol:eje:1-correos}{{A.1}{40}{Tema 1}{section*.6}{}}
\newlabel{sol:eje:1-correos@cref}{{[subappendix][1][2147483647,1]A.1}{[1][40][]40}}
\newlabel{sol:eje:1-piensos}{{A.1}{41}{Tema 1}{section*.7}{}}
\newlabel{sol:eje:1-piensos@cref}{{[subappendix][1][2147483647,1]A.1}{[1][41][]41}}
\newlabel{sol:eje:1-minimax}{{A.1}{41}{Tema 1}{section*.8}{}}
\newlabel{sol:eje:1-minimax@cref}{{[subappendix][1][2147483647,1]A.1}{[1][41][]41}}
\newlabel{sol:eje:1-desviaciones}{{A.1}{42}{Tema 1}{section*.9}{}}
\newlabel{sol:eje:1-desviaciones@cref}{{[subappendix][1][2147483647,1]A.1}{[1][41][]42}}
\newlabel{sol:eje:1-diseno}{{A.1}{42}{Tema 1}{section*.10}{}}
\newlabel{sol:eje:1-diseno@cref}{{[subappendix][1][2147483647,1]A.1}{[1][42][]42}}
\newlabel{sol:eje:2-holgura}{{A.2}{43}{Tema 2}{section*.11}{}}
\newlabel{sol:eje:2-holgura@cref}{{[subappendix][2][2147483647,1]A.2}{[1][42][]43}}
\newlabel{sol:eje:2-dual}{{A.2}{43}{Tema 2}{section*.12}{}}
\newlabel{sol:eje:2-dual@cref}{{[subappendix][2][2147483647,1]A.2}{[1][43][]43}}
\newlabel{sol:eje:2-outdoors}{{A.2}{43}{Tema 2}{section*.13}{}}
\newlabel{sol:eje:2-outdoors@cref}{{[subappendix][2][2147483647,1]A.2}{[1][43][]43}}
\newlabel{sol:eje:2-bb}{{A.2}{44}{Tema 2}{section*.14}{}}
\newlabel{sol:eje:2-bb@cref}{{[subappendix][2][2147483647,1]A.2}{[1][44][]44}}
\newlabel{sol:eje:2-mochila}{{A.2}{44}{Tema 2}{section*.15}{}}
\newlabel{sol:eje:2-mochila@cref}{{[subappendix][2][2147483647,1]A.2}{[1][44][]44}}
\newlabel{sol:eje:2-logica}{{A.2}{45}{Tema 2}{section*.16}{}}
\newlabel{sol:eje:2-logica@cref}{{[subappendix][2][2147483647,1]A.2}{[1][45][]45}}
\newlabel{sol:eje:2-bigm}{{A.2}{45}{Tema 2}{section*.17}{}}
\newlabel{sol:eje:2-bigm@cref}{{[subappendix][2][2147483647,1]A.2}{[1][45][]45}}
\newlabel{sol:eje:2-concava}{{A.2}{45}{Tema 2}{section*.18}{}}
\newlabel{sol:eje:2-concava@cref}{{[subappendix][2][2147483647,1]A.2}{[1][45][]45}}
\newlabel{sol:eje:3-contar}{{A.3}{46}{Tema 3}{section*.19}{}}
\newlabel{sol:eje:3-contar@cref}{{[subappendix][3][2147483647,1]A.3}{[1][46][]46}}
\newlabel{sol:eje:3-ruta}{{A.3}{46}{Tema 3}{section*.20}{}}
\newlabel{sol:eje:3-ruta@cref}{{[subappendix][3][2147483647,1]A.3}{[1][46][]46}}
\newlabel{sol:eje:3-2opt}{{A.3}{47}{Tema 3}{section*.21}{}}
\newlabel{sol:eje:3-2opt@cref}{{[subappendix][3][2147483647,1]A.3}{[1][46][]47}}
\newlabel{sol:eje:3-cuñas}{{A.3}{47}{Tema 3}{section*.22}{}}
\newlabel{sol:eje:3-cuñas@cref}{{[subappendix][3][2147483647,1]A.3}{[1][47][]47}}
\newlabel{sol:eje:3-johnson}{{A.3}{47}{Tema 3}{section*.23}{}}
\newlabel{sol:eje:3-johnson@cref}{{[subappendix][3][2147483647,1]A.3}{[1][47][]47}}
\newlabel{sol:eje:3-bigm}{{A.3}{48}{Tema 3}{section*.24}{}}
\newlabel{sol:eje:3-bigm@cref}{{[subappendix][3][2147483647,1]A.3}{[1][47][]48}}
\newlabel{sol:eje:3-cota}{{A.3}{48}{Tema 3}{section*.25}{}}
\newlabel{sol:eje:3-cota@cref}{{[subappendix][3][2147483647,1]A.3}{[1][48][]48}}
\newlabel{sol:eje:4-dcp}{{A.4}{49}{Tema 4}{section*.26}{}}
\newlabel{sol:eje:4-dcp@cref}{{[subappendix][4][2147483647,1]A.4}{[1][49][]49}}
\newlabel{sol:eje:4-kkt}{{A.4}{49}{Tema 4}{section*.27}{}}
\newlabel{sol:eje:4-kkt@cref}{{[subappendix][4][2147483647,1]A.4}{[1][49][]49}}
\newlabel{sol:eje:4-despacho2}{{A.4}{49}{Tema 4}{section*.28}{}}
\newlabel{sol:eje:4-despacho2@cref}{{[subappendix][4][2147483647,1]A.4}{[1][49][]49}}
\newlabel{sol:eje:4-agua}{{A.4}{50}{Tema 4}{section*.29}{}}
\newlabel{sol:eje:4-agua@cref}{{[subappendix][4][2147483647,1]A.4}{[1][50][]50}}
\newlabel{sol:eje:4-sensib}{{A.4}{50}{Tema 4}{section*.30}{}}
\newlabel{sol:eje:4-sensib@cref}{{[subappendix][4][2147483647,1]A.4}{[1][50][]50}}
\newlabel{sol:eje:4-dual}{{A.4}{51}{Tema 4}{section*.31}{}}
\newlabel{sol:eje:4-dual@cref}{{[subappendix][4][2147483647,1]A.4}{[1][51][]51}}
\newlabel{sol:eje:4-primer-orden}{{A.4}{51}{Tema 4}{section*.32}{}}
\newlabel{sol:eje:4-primer-orden@cref}{{[subappendix][4][2147483647,1]A.4}{[1][51][]51}}
\newlabel{sol:eje:5-aceptacion}{{A.5}{52}{Tema 5}{section*.33}{}}
\newlabel{sol:eje:5-aceptacion@cref}{{[subappendix][5][2147483647,1]A.5}{[1][51][]52}}
\newlabel{sol:eje:5-frentes}{{A.5}{52}{Tema 5}{section*.34}{}}
\newlabel{sol:eje:5-frentes@cref}{{[subappendix][5][2147483647,1]A.5}{[1][52][]52}}
\newlabel{sol:eje:5-crowding}{{A.5}{52}{Tema 5}{section*.35}{}}
\newlabel{sol:eje:5-crowding@cref}{{[subappendix][5][2147483647,1]A.5}{[1][52][]52}}
\newlabel{sol:eje:5-apoyado}{{A.5}{52}{Tema 5}{section*.36}{}}
\newlabel{sol:eje:5-apoyado@cref}{{[subappendix][5][2147483647,1]A.5}{[1][52][]52}}
\newlabel{sol:eje:5-ox}{{A.5}{53}{Tema 5}{section*.37}{}}
\newlabel{sol:eje:5-ox@cref}{{[subappendix][5][2147483647,1]A.5}{[1][53][]53}}
\newlabel{sol:eje:5-pesos}{{A.5}{53}{Tema 5}{section*.38}{}}
\newlabel{sol:eje:5-pesos@cref}{{[subappendix][5][2147483647,1]A.5}{[1][53][]53}}
\newlabel{sol:eje:5-eps}{{A.5}{53}{Tema 5}{section*.39}{}}
\newlabel{sol:eje:5-eps@cref}{{[subappendix][5][2147483647,1]A.5}{[1][53][]53}}
\newlabel{sol:eje:5-metropolis}{{A.5}{54}{Tema 5}{section*.40}{}}
\newlabel{sol:eje:5-metropolis@cref}{{[subappendix][5][2147483647,1]A.5}{[1][53][]54}}
\newlabel{sol:eje:6-fugas}{{A.6}{54}{Tema 6}{section*.41}{}}
\newlabel{sol:eje:6-fugas@cref}{{[subappendix][6][2147483647,1]A.6}{[1][54][]54}}
\newlabel{sol:eje:6-newsvendor}{{A.6}{55}{Tema 6}{section*.42}{}}
\newlabel{sol:eje:6-newsvendor@cref}{{[subappendix][6][2147483647,1]A.6}{[1][54][]55}}
\newlabel{sol:eje:6-spo}{{A.6}{55}{Tema 6}{section*.43}{}}
\newlabel{sol:eje:6-spo@cref}{{[subappendix][6][2147483647,1]A.6}{[1][55][]55}}
\newlabel{sol:eje:6-vss}{{A.6}{55}{Tema 6}{section*.44}{}}
\newlabel{sol:eje:6-vss@cref}{{[subappendix][6][2147483647,1]A.6}{[1][55][]55}}
\newlabel{sol:eje:6-robusta}{{A.6}{56}{Tema 6}{section*.45}{}}
\newlabel{sol:eje:6-robusta@cref}{{[subappendix][6][2147483647,1]A.6}{[1][56][]56}}
\newlabel{sol:eje:6-saa}{{A.6}{56}{Tema 6}{section*.46}{}}
\newlabel{sol:eje:6-saa@cref}{{[subappendix][6][2147483647,1]A.6}{[1][56][]56}}
\newlabel{sol:eje:6-fractil}{{A.6}{56}{Tema 6}{section*.47}{}}
\newlabel{sol:eje:6-fractil@cref}{{[subappendix][6][2147483647,1]A.6}{[1][56][]56}}
\newlabel{sol:eje:6-maldicion}{{A.6}{57}{Tema 6}{section*.48}{}}
\newlabel{sol:eje:6-maldicion@cref}{{[subappendix][6][2147483647,1]A.6}{[1][57][]57}}
\makeatother

\begin{document}
\frontmatter
\begin{titlepage}
\centering
\vspace*{2.5cm}
{\large\scshape Anatomía de un LLM}\par
\vspace{0.4cm}
{\normalsize Fundamentos matemáticos de los modelos de lenguaje}\par
\vspace{2.2cm}
\rule{\textwidth}{0.8pt}\par\vspace{0.5cm}
{\small\scshape Fase VI \textperiodcentered\ Diseño y construcción de optimizadores con IA}\par\vspace{0.3cm}
{\Huge\bfseries Módulo 24\par}\vspace{0.3cm}
{\LARGE Formular y resolver problemas de optimización\par}
\vspace{0.2cm}
{\large Programación lineal y entera, restricciones, convexidad, metaheurísticas y decisiones con predicciones\par}
\vspace{0.5cm}\rule{\textwidth}{0.4pt}\par
\vspace{1.6cm}
{\large Resumen del módulo}\par
\vfill
{\small Edición de octubre de 2026}\par
\end{titlepage}

\chapter*{Sobre este resumen}
\addcontentsline{toc}{chapter}{Sobre este resumen}
\markboth{Sobre este resumen}{}

Este documento reúne lo indispensable de los apuntes del Módulo 24: el texto expositivo, todas las definiciones y todos los resultados (teoremas, proposiciones, lemas y corolarios), las tablas y las observaciones de contenido. Se han suprimido las demostraciones, los ejemplos, las lecturas, las figuras, los ejercicios y sus soluciones. La numeración de temas, secciones, enunciados y ecuaciones es la misma que en los apuntes, de modo que cualquier referencia puede buscarse allí; las que remiten a elementos suprimidos (un ejemplo, una figura, una ecuación de una demostración) conservan su número en los apuntes.

\section*{Notación}
\begin{center}\small
\begin{tabularx}{\textwidth}{@{}l>{\raggedright\arraybackslash}X@{}}
\toprule
Símbolo & Significado \\
\midrule
$x\in\R^{n}$, $\mathcal{X}$ & vector de decisión y conjunto factible \\
$f$, $g_i$, $h_j$ & función objetivo, restricciones de desigualdad ($g_i\le 0$) y de igualdad ($h_j=0$) \\
$\theta$, $\hat\theta$ & parámetros (datos) del problema y su estimación o predicción \\
$c,A,b$ & datos de un programa lineal $\max\{c\T x: Ax\le b,\ x\ge 0\}$ \\
$y$, $y^{*}$ & variables duales; precios sombra en el óptimo \\
$z^{*}$, $z_{\mathrm{LP}}$ & valor óptimo de un problema y de su relajación lineal \\
$M$ & constante de las restricciones big-M \\
$\Lag(x,\lambda,\nu)$ & lagrangiana con multiplicadores $\lambda\ge 0$ (desigualdades) y $\nu$ (igualdades) \\
$T$ & temperatura del recocido simulado \\
$\succ$ & dominancia de Pareto ($u\succ v$: $u$ domina a $v$) \\
$\xi$, $\Xi$ & parámetro incierto y conjunto de incertidumbre \\
$c_u$, $c_o$ & coste unitario por defecto y por exceso (vendedor de periódicos) \\
$F$, $F^{-1}$ & función de distribución y función cuantil \\
hL, MW, MWh, h & hectolitros, megavatios, megavatios hora, horas \\
\bottomrule
\end{tabularx}
\end{center}


\tableofcontents

\mainmatter
% =====================================================================
\chapter{Del problema real al modelo matemático}
\label{tema:1}
% =====================================================================

Antes de elegir un algoritmo hay que saber qué se decide, qué se quiere mejorar y qué no se puede violar. Este tema fija el vocabulario de la modelización siguiendo el capítulo~1 de Bertsimas y Tsitsiklis~\cite{bt}, la clase~1 («Formulations») del curso 6.251J del MIT~\cite{mit6251} y la clase~2 («Formulations of linear and non-linear programs») del curso 15.053~\cite{mit15053}; describe la arquitectura de un sistema que predice y decide, con dos aplicaciones industriales publicadas~\cite{deepmind,orion}, y traza el mapa de métodos que recorre la fase, con el capítulo~1 de Kochenderfer y Wheeler~\cite{kw} como guía.

\section{Elementos de un modelo de optimización}

Un problema real llega como una descripción verbal: «queremos producir lo más rentable posible sin parar la línea y sin dejar a ningún cliente sin pedido». Modelizar consiste en separar lo que se elige de lo que se sabe y en escribir con precisión lo que se quiere y lo que se exige.

\begin{definicion}[Problema de optimización]\label{def:prob-opt}
Un \emph{problema de optimización} es
\[
\min_{x\in X}\ f(x)\qquad \st g_i(x)\le 0\ \ (i=1,\dots,m),\qquad h_j(x)=0\ \ (j=1,\dots,p),
\]
donde $x$ es el vector de \emph{variables de decisión}, $X$ un conjunto básico (por ejemplo $\R^{n}$, $\R^{n}_{\ge 0}$ o $\Z^{n_1}\times\R^{n_2}$), $f$ la \emph{función objetivo} y $g_i,h_j$ las \emph{restricciones}. El \emph{conjunto factible} es $\mathcal{X}=\{x\in X: g_i(x)\le 0,\ h_j(x)=0\}$. Un punto $x^{*}\in\mathcal{X}$ es \emph{óptimo (global)} si $f(x^{*})\le f(x)$ para todo $x\in\mathcal{X}$, y \emph{óptimo local} si la desigualdad se cumple en $\mathcal{X}\cap B(x^{*},r)$ para algún $r>0$. El problema es \emph{infactible} si $\mathcal{X}=\varnothing$ y \emph{no acotado} si $\inf_{\mathcal{X}}f=-\infty$.
\end{definicion}

\begin{definicion}[Datos, decisiones e indicadores]\label{def:elementos}
En un modelo se distinguen:
\begin{enumerate}[label=(\alph*)]
  \item los \emph{parámetros} o datos $\theta$ (precios, capacidades, demandas, tiempos), que el decisor no controla y que pueden ser inciertos o predichos;
  \item las \emph{variables de decisión} $x$, que el decisor fija;
  \item las \emph{restricciones duras}, que ninguna decisión puede violar (seguridad, física, contratos), y las \emph{restricciones blandas}, que pueden violarse con una penalización en el objetivo;
  \item el \emph{objetivo} $f(x;\theta)$, un único escalar que se optimiza;
  \item los \emph{indicadores clave} (KPIs, \emph{key performance indicators}): magnitudes que se miden para vigilar el sistema (utilización, nivel de servicio, emisiones, tiempo de cálculo) y que no tienen por qué coincidir con el objetivo.
\end{enumerate}
\end{definicion}

\begin{observacion}
El objetivo es una decisión de diseño: si se optimiza el coste y se vigila el nivel de servicio, conviene comprobar que el óptimo del coste no hunde el servicio. Cuando varios KPIs son igualmente importantes, el problema es multiobjetivo (\cref{tema:5}). Optimizar un indicador que solo aproxima lo que se quiere conduce a la ley de Goodhart que el módulo~21 estudió como sobreoptimización de la recompensa.
\end{observacion}

Muchos modelos que no parecen lineales lo son tras una reformulación. Las tres siguientes son las más útiles.

\begin{proposicion}[Reformulaciones elementales]\label{prop:reform}
\begin{enumerate}
  \item Maximizar $f$ equivale a minimizar $-f$: $\max_{\mathcal{X}}f=-\min_{\mathcal{X}}(-f)$, con los mismos puntos óptimos.
  \item Una igualdad $h(x)=0$ equivale a las dos desigualdades $h(x)\le 0$ y $-h(x)\le 0$; una desigualdad $a\T x\le b$ equivale a $a\T x+s=b$ con una \emph{holgura} $s\ge 0$.
  \item Una variable libre $x_j\in\R$ puede sustituirse por $x_j=x_j^{+}-x_j^{-}$ con $x_j^{\pm}\ge 0$.
\end{enumerate}
\end{proposicion}

\begin{proposicion}[Epígrafe de un máximo]\label{prop:minimax}
Sean $f_1,\dots,f_k$ funciones sobre $\mathcal{X}$. Los problemas
\[
\text{(P)}\ \ \min_{x\in\mathcal{X}}\ \max_{i}f_i(x)\qquad\text{y}\qquad \text{(Q)}\ \ \min_{x\in\mathcal{X},\,t\in\R}\ t\ \ \st f_i(x)\le t\ \ (i=1,\dots,k)
\]
tienen el mismo valor óptimo, y $x^{*}$ es óptimo de (P) si y solo si $(x^{*},\max_i f_i(x^{*}))$ es óptimo de (Q). Si las $f_i$ son afines y $\mathcal{X}$ es un poliedro, (Q) es un programa lineal.
\end{proposicion}

\begin{proposicion}[Suma de valores absolutos]\label{prop:abs}
El problema $\min_{x\in\mathcal{X}}\sum_{i=1}^{k}\abs{a_i\T x-b_i}$ equivale a
\[
\min_{x\in\mathcal{X},\,e\in\R^{k}}\ \sum_{i=1}^{k}e_i\qquad\st -e_i\le a_i\T x-b_i\le e_i\quad(i=1,\dots,k).
\]
\end{proposicion}

\addtocounter{definicion}{1}%
\addtocounter{table}{1}%

\addtocounter{figure}{1}%

\section{Anatomía de un optimizador con IA}

En la industria el modelo matemático es una pieza de un sistema mayor que se ejecuta de forma periódica (cada semana, cada hora, cada cinco minutos). La \cref{fig:anatomia} muestra sus componentes.

\addtocounter{figure}{1}%

\begin{definicion}[Optimizador con IA]\label{def:optimizador-ia}
Un \emph{optimizador con IA} es un procedimiento que, a partir de observaciones $w$ (contexto: fecha, temperatura, estado de la planta), produce una decisión
\[
\hat x(w)\in\argmin_{x\in\mathcal{X}(\hat\theta)}\ f\bigl(x;\hat\theta\bigr),\qquad \hat\theta=m_\phi(w),
\]
donde $m_\phi$ es un modelo aprendido de datos con parámetros $\phi$ que predice los parámetros inciertos del problema (o el propio valor de una decisión), y el conjunto factible incluye las restricciones de seguridad. La calidad del sistema se mide por el coste real $f(\hat x(w);\theta)$ con el $\theta$ verdadero, no por el error de predicción.
\end{definicion}

Los componentes de la \cref{fig:anatomia} son:
\begin{enumerate}[label=\arabic*.]
  \item \emph{Datos}: históricos, tiempo real y datos maestros (capacidades, recetas, tarifas). Su calidad limita todo lo demás.
  \item \emph{Modelo predictivo}: estima demandas, tiempos, rendimientos o el estado futuro del sistema (regresión, árboles, redes; módulos~8--11).
  \item \emph{Modelo de optimización}: la formulación de la \cref{def:prob-opt}, con los parámetros predichos y las restricciones duras.
  \item \emph{Resolvedor}: el algoritmo elegido según la estructura (\cref{sec:mapa}).
  \item \emph{Decisión con salvaguardas}: comprobaciones independientes del modelo (límites físicos, cambios máximos entre periodos, plan de respaldo) antes de ejecutar.
  \item \emph{Monitorización}: KPIs, error de predicción, deriva de los datos, tiempo de cálculo y tasa de intervenciones humanas; disparan el reentrenamiento o la revisión del modelo.
\end{enumerate}

\addtocounter{definicion}{1}%

\section{Mapa de métodos de la fase}\label{sec:mapa}

La estructura matemática del problema, no la moda, decide el método. El \cref{tab:mapa} resume la elección; la \cref{fig:arbol} la presenta como árbol de decisión.

\begin{table}[h]
\centering\small
\caption{Mapa de métodos: estructura del problema, familia y herramienta típica.}\label{tab:mapa}
\begin{tabularx}{\textwidth}{@{}>{\raggedright\arraybackslash}p{3.5cm}>{\raggedright\arraybackslash}p{3.1cm}>{\raggedright\arraybackslash}X@{}}
\toprule
Estructura & Familia & Métodos y herramientas \\
\midrule
objetivo y restricciones lineales, variables continuas & PL (LP) & símplex, punto interior; HiGHS (\cref{tema:2}) \\
lineal con decisiones discretas (sí/no, lotes, turnos) & PLEM (MILP) & ramificación y acotación con cortes; HiGHS, SCIP, Gurobi (\cref{tema:2}) \\
secuenciación, rutas, asignación con lógica compleja & combinatoria, programación por restricciones & CP-SAT, OR-Tools Routing, búsqueda local (\cref{tema:3}) \\
convexo y diferenciable (cuadrático, cónico) & optimización convexa & punto interior; CVXPY con Clarabel u OSQP (\cref{tema:4}) \\
no convexo y diferenciable & programación no lineal & SQP, punto interior, multiarranque; SciPy (\cref{tema:4}) \\
caja negra barata, sin derivadas & metaheurísticas & recocido, genéticos, CMA-ES, NSGA-II (\cref{tema:5}) \\
caja negra cara (simulación, experimento) & optimización bayesiana & modelo sustituto gaussiano y función de adquisición~\cite{garnett,frazier} \\
decisiones secuenciales con incertidumbre & control y RL & control predictivo, aprendizaje por refuerzo (módulo~21)~\cite{sb} \\
parámetros inciertos que se predicen & predecir y optimizar, estocástica, robusta & SAA, recurso, conjuntos de incertidumbre (\cref{tema:6}) \\
espacio de programas o heurísticas & búsqueda con LLM y evaluador & FunSearch~\cite{funsearch}, AlphaEvolve~\cite{alphaevolve}, OPRO~\cite{opro} \\
\bottomrule
\end{tabularx}
\end{table}

\addtocounter{figure}{1}%

\begin{observacion}[Modelos de lenguaje en la optimización]\label{obs:llm}
Los LLM intervienen en la optimización de dos formas distintas. Como \emph{asistentes de formulación}, traducen un enunciado en lenguaje natural a un modelo para un resolvedor clásico; OptiMUS~\cite{optimus} es un ejemplo publicado para PL y PLEM, y su salida debe verificarse como cualquier código. Como \emph{generadores de candidatos} dentro de un bucle evolutivo con un evaluador automático, proponen programas o soluciones que el evaluador puntúa: FunSearch~\cite{funsearch} encontró así construcciones matemáticas y heurísticas de empaquetado, AlphaEvolve~\cite{alphaevolve} extendió la idea a algoritmos completos y OPRO~\cite{opro} usa el propio LLM como optimizador de instrucciones. En ambos casos la garantía procede del evaluador o del resolvedor, no del modelo de lenguaje; los módulos posteriores de la fase desarrollan este enfoque.
\end{observacion}

\begin{table}[h]
\centering\small
\caption{Tecnologías del tema~1: lenguajes de modelización y resolvedores generales.}\label{tab:tec1}
\begin{tabularx}{\textwidth}{@{}l>{\raggedright\arraybackslash}X>{\raggedright\arraybackslash}p{2.2cm}>{\raggedright\arraybackslash}X@{}}
\toprule
Herramienta & Para qué sirve & Licencia o tipo & Cuándo elegirla \\
\midrule
SciPy (\texttt{optimize}) & PL y PLEM con HiGHS, mínimos cuadrados, no lineal local & abierta (BSD-3) & prototipos y modelos pequeños sin dependencias extra \\
PuLP~\cite{pulp} & modelizar PL/PLEM y enviarlos a CBC, HiGHS, SCIP, Gurobi\dots & abierta (MIT) & modelos lineales legibles en Python \\
Pyomo & modelización algebraica general (lineal y no lineal) & abierta (BSD-3) & modelos grandes, con estructura por índices y varios resolvedores \\
OR-Tools~\cite{ortools-cpsat} & CP-SAT, rutas, flujos, PL (GLOP) & abierta (Apache~2.0) & combinatoria, programación y rutas \\
CVXPY~\cite{cvxpy} & optimización convexa disciplinada & abierta (Apache~2.0) & problemas convexos con verificación automática \\
Gurobi, CPLEX & resolvedores comerciales de PL/PLEM/cuadráticos & comercial (licencia académica gratuita) & PLEM grandes en producción con exigencia de tiempo \\
\bottomrule
\end{tabularx}
\end{table}

\addtocounter{definicion}{1}%


% =====================================================================
\chapter{Programación lineal y entera mixta}
\label{tema:2}
% =====================================================================

La programación lineal y la entera mixta resuelven la mayor parte de los problemas de planificación de la industria: producción, mezclas, transporte, turnos, inversiones. Este tema sigue los capítulos~2 a~5, 10 y~11 de Bertsimas y Tsitsiklis~\cite{bt}, que son las lecturas del curso 6.251J del MIT~\cite{mit6251}, y las clases~6 («Sensitivity analysis and shadow prices»), 10--12 (programación entera y ramificación y acotación) del curso 15.053~\cite{mit15053}, con los capítulos~3 («Sensitivity Analysis») y~9 («Integer Programming») de Bradley, Hax y Magnanti~\cite{amp}. La ramificación y acotación procede de Land y Doig~\cite{landdoig}; el resolvedor de referencia es HiGHS~\cite{highs,hh}.

\section{Programas lineales y vértices}

\begin{definicion}[Programa lineal]\label{def:pl}
Un \emph{programa lineal} (PL) en forma canónica es
\[
z^{*}=\max\ c\T x\qquad\st Ax\le b,\quad x\ge 0,
\]
con $A\in\R^{m\times n}$, $b\in\R^{m}$ y $c\in\R^{n}$. Su conjunto factible $\mathcal{P}=\{x: Ax\le b,\ x\ge 0\}$ es un \emph{poliedro}. Un punto $v\in\mathcal{P}$ es un \emph{vértice} si no es combinación convexa $\lambda u+(1-\lambda)w$ de dos puntos distintos $u,w\in\mathcal{P}$ con $\lambda\in(0,1)$; equivalentemente, si en $v$ hay $n$ restricciones activas linealmente independientes~\cite{bt}.
\end{definicion}

\begin{teorema}[Teorema fundamental de la programación lineal]\label{teo:vertice}
Si $\mathcal{P}$ es no vacío y acotado, el PL tiene un óptimo que es un vértice de $\mathcal{P}$.
\end{teorema}

El \emph{método símplex} (Dantzig) explota el teorema: recorre vértices adyacentes mejorando el objetivo hasta que ninguna arista mejora. Los \emph{métodos de punto interior} avanzan por el interior del poliedro y tienen complejidad polinómica. HiGHS implementa ambos (su símplex dual revisado paralelo se describe en~\cite{hh}) y un resolvedor de PLEM; SciPy lo usa en \texttt{linprog} y \texttt{milp}~\cite{highs}.

\section{Dualidad y precios sombra}

\begin{definicion}[Problema dual]\label{def:dual}
El \emph{dual} del PL de la \cref{def:pl} es
\[
w^{*}=\min\ b\T y\qquad\st A\T y\ge c,\quad y\ge 0.
\]
Cada variable dual $y_i$ corresponde a la restricción $i$ del primal.
\end{definicion}

\begin{teorema}[Dualidad débil]\label{teo:dual-debil}
Si $x$ es factible para el primal e $y$ para el dual, entonces $c\T x\le b\T y$. En consecuencia $z^{*}\le w^{*}$, y si $c\T x=b\T y$ ambos son óptimos.
\end{teorema}

\begin{teorema}[Dualidad fuerte]\label{teo:dual-fuerte}
Si el primal tiene óptimo, también lo tiene el dual y $z^{*}=w^{*}$ (Bertsimas y Tsitsiklis~\cite{bt}, cap.~4).
\end{teorema}

\begin{corolario}[Holgura complementaria]\label{cor:holgura}
Sean $x$ e $y$ factibles. Ambos son óptimos si y solo si
\[
y_i\,\bigl(b_i-(Ax)_i\bigr)=0\ \ \text{para todo }i\qquad\text{y}\qquad x_j\,\bigl((A\T y)_j-c_j\bigr)=0\ \ \text{para todo }j.
\]
\end{corolario}

La holgura complementaria tiene una lectura económica directa: un recurso que sobra (restricción no activa) tiene precio cero, y un producto que no se fabrica no cubre con su margen el valor de los recursos que consumiría.

\begin{proposicion}[Precios sombra]\label{prop:sombra}
Sea $z^{*}(b)$ el valor óptimo en función del término independiente. La función $z^{*}$ es cóncava en el conjunto de $b$ con primal factible y acotado, y si $y^{*}$ es una solución dual óptima para $b$, para todo $b'$ de ese conjunto
\[
z^{*}(b')\le z^{*}(b)+(y^{*})\T(b'-b).
\]
Si el dual óptimo es único, $z^{*}$ es diferenciable en $b$ con $\partial z^{*}/\partial b_i=y_i^{*}$: el \emph{precio sombra} $y_i^{*}$ es lo que aumenta el valor óptimo por unidad adicional del recurso $i$, mientras la base óptima no cambia.
\end{proposicion}

\addtocounter{definicion}{1}%

\section{Programación entera mixta}

\begin{definicion}[PLEM y relajación lineal]\label{def:plem}
Un \emph{programa lineal entero mixto} (PLEM, en inglés MILP) es un PL en el que un subconjunto $J$ de variables debe ser entero: $\max\{c\T x: Ax\le b,\ x\ge 0,\ x_j\in\Z\ (j\in J)\}$. Su \emph{relajación lineal} es el mismo problema sin la condición de integralidad; su valor se denota $z_{\mathrm{LP}}$.
\end{definicion}

\begin{proposicion}[La relajación acota]\label{prop:relajacion}
En un problema de maximización, $z^{*}\le z_{\mathrm{LP}}$; en uno de minimización, $z^{*}\ge z_{\mathrm{LP}}$. Si la solución de la relajación cumple la integralidad, es óptima para el PLEM.
\end{proposicion}

La diferencia relativa entre el mejor valor entero conocido (el \emph{incumbente}) y la mejor cota se llama \emph{gap}; los resolvedores se detienen cuando es menor que una tolerancia (por defecto, del orden de $10^{-4}$).

Las variables binarias $y\in\{0,1\}$ permiten modelizar decisiones lógicas con restricciones lineales:
\begin{enumerate}[label=(\alph*)]
  \item \emph{coste fijo o de preparación}: $x\le My$ con coste $fy$ en el objetivo; si se produce ($x>0$), se paga $f$;
  \item \emph{lote mínimo}: $Ly\le x\le My$; o no se produce, o se producen al menos $L$ unidades;
  \item \emph{implicación} «si se elige $A$, se elige $B$»: $y_A\le y_B$; \emph{exclusión}: $y_A+y_B\le 1$; \emph{al menos $k$ de $n$}: $\sum_i y_i\ge k$;
  \item \emph{disyunción} «se cumple $a\T x\le b$ o $d\T x\le e$»: $a\T x\le b+M\delta$, $d\T x\le e+M(1-\delta)$, $\delta\in\{0,1\}$.
\end{enumerate}

\begin{proposicion}[Validez y fuerza de una cota big-M]\label{prop:bigm}
Sea $\mathcal{X}$ el conjunto factible sin la restricción $x\le My$ y sea $U=\sup\{x: x\in\mathcal{X}\}$.
\begin{enumerate}
  \item La restricción $x\le My$, con $y\in\{0,1\}$, no elimina ninguna solución con $y=1$ si y solo si $M\ge U$.
  \item Si $U\le M\le M'$, la relajación lineal con $M$ está contenida en la relajación con $M'$; la menor cota válida, $M=U$, da la relajación más ajustada.
\end{enumerate}
\end{proposicion}

Con una $M$ grande la relajación permite valores de $y$ diminutos, por ejemplo $y=x/M$, y paga solo una fracción del coste fijo: la cota es mala y el árbol de búsqueda crece. Una $M$ desmesurada causa además errores numéricos, porque los resolvedores trabajan con tolerancias del orden de $10^{-6}$ a $10^{-9}$.

\section{Ramificación y acotación}

La \emph{ramificación y acotación} (\emph{branch and bound})~\cite{landdoig} explora un árbol de subproblemas. En cada nodo se resuelve la relajación lineal con las cotas acumuladas. Para un problema de maximización:
\begin{enumerate}[label=\arabic*.]
  \item si la relajación es infactible, el nodo se \emph{poda};
  \item si su valor no supera el del incumbente, se poda \emph{por cota}: ningún descendiente puede mejorar (\cref{prop:relajacion});
  \item si la solución es entera, se actualiza el incumbente y se poda;
  \item si no, se elige una variable $x_j$ con valor fraccionario $x_j^{*}$ y se crean dos hijos con $x_j\le\lfloor x_j^{*}\rfloor$ y $x_j\ge\lceil x_j^{*}\rceil$.
\end{enumerate}
Si todas las variables enteras están acotadas, el árbol es finito y el incumbente final es óptimo: cada nodo podado lo es con una justificación válida y las ramas cubren todos los enteros factibles. Si los coeficientes del objetivo son enteros, el valor de cualquier solución entera también lo es, y un nodo puede podarse cuando $\lfloor z_{\mathrm{LP}}\rfloor$ no supera al incumbente. Los resolvedores modernos añaden preproceso, planos de corte (desigualdades válidas que eliminan soluciones fraccionarias; clase~13 de~\cite{mit15053}), heurísticas primales y reglas de ramificación elaboradas.

\addtocounter{definicion}{1}%

\begin{table}[h]
\centering\small
\caption{Tecnologías del tema~2: programación lineal y entera mixta.}\label{tab:tec2}
\begin{tabularx}{\textwidth}{@{}l>{\raggedright\arraybackslash}X>{\raggedright\arraybackslash}p{2.3cm}>{\raggedright\arraybackslash}X@{}}
\toprule
Herramienta & Para qué sirve & Licencia o tipo & Cuándo elegirla \\
\midrule
HiGHS~\cite{highs} & PL (símplex dual, punto interior), PLEM, cuadrático & abierta (MIT) & resolvedor abierto de referencia; incluido en SciPy \\
SciPy \texttt{linprog}/\texttt{milp} & interfaz matricial a HiGHS, con duales & abierta (BSD-3) & modelos pequeños y docencia \\
PuLP~\cite{pulp} & modelización por nombres; CBC, HiGHS, SCIP\dots & abierta (MIT) & modelos de planificación legibles \\
SCIP & PLEM y no lineal entero, con cortes y ramificación configurables & abierta (Apache~2.0 desde la 8.0.3) & investigación y PLEM difíciles sin licencia comercial \\
Gurobi, CPLEX & PL/PLEM a gran escala & comercial & producción con exigencias de tiempo y soporte \\
OR-Tools (GLOP, MathOpt) & PL y envoltorio de resolvedores & abierta (Apache~2.0) & si ya se usa OR-Tools para rutas o CP \\
\bottomrule
\end{tabularx}
\end{table}

\addtocounter{definicion}{1}%


% =====================================================================
\chapter{Problemas combinatorios y programación por restricciones}
\label{tema:3}
% =====================================================================

Rutas, secuencias y calendarios tienen un número de soluciones que crece de forma factorial con el tamaño. Este tema estudia cómo formularlos y qué herramientas los resuelven en la práctica. Sigue la clase~17 («Traveling salesman problem») del curso 15.053 del MIT~\cite{mit15053} para la complejidad y las heurísticas, el capítulo~10 de Bertsimas y Tsitsiklis~\cite{bt} para las formulaciones enteras, el artículo de Solomon~\cite{solomon} para las rutas con ventanas horarias, el de Johnson~\cite{johnson} para el taller con dos máquinas y la documentación de OR-Tools~\cite{ortools-vrptw,ortools-js,ortools-cpsat} para la programación por restricciones.

\section{Explosión combinatoria}

\begin{proposicion}[Número de circuitos]\label{prop:circuitos}
Con $n\ge 3$ ciudades y distancias simétricas hay $(n-1)!/2$ circuitos hamiltonianos distintos.
\end{proposicion}

Con $n=20$ hay $19!/2\approx 6{,}1\cdot 10^{16}$ circuitos. El problema del viajante (TSP) es NP-difícil: no se conoce un algoritmo polinómico, y encontrarlo daría uno para todos los problemas NP-completos~\cite{mit15053}. En la práctica se combinan tres enfoques:
\begin{enumerate}[label=(\alph*)]
  \item \emph{métodos exactos}: PLEM con ramificación, cortes y generación de columnas, o programación dinámica para instancias pequeñas; dan el óptimo y una cota;
  \item \emph{heurísticas constructivas}: vecino más próximo, inserción, ahorros; dan una solución rápida;
  \item \emph{búsqueda local y metaheurísticas}: intercambios 2-opt y 3-opt, recolocación de clientes, búsqueda local guiada, recocido (\cref{tema:5}).
\end{enumerate}

\begin{definicion}[Movimiento 2-opt]\label{def:2opt}
Dado un circuito $(\dots,a,b,\dots,c,d,\dots)$, el movimiento 2-opt elimina las aristas $(a,b)$ y $(c,d)$ y añade $(a,c)$ y $(b,d)$, invirtiendo el tramo de $b$ a $c$. La longitud cambia en $\Delta=d(a,c)+d(b,d)-d(a,b)-d(c,d)$; el movimiento se aplica si $\Delta<0$. Un circuito es \emph{2-óptimo} si ningún movimiento 2-opt lo mejora.
\end{definicion}

En el plano euclídeo, un circuito con dos aristas que se cruzan nunca es 2-óptimo: deshacer el cruce lo acorta, por la desigualdad triangular aplicada a los dos triángulos que forman las diagonales (\cref{eje:3-2opt}).

\section{Rutas de vehículos con ventanas horarias}

\begin{definicion}[VRPTW]\label{def:vrptw}
En el \emph{problema de rutas de vehículos con ventanas horarias} (VRPTW) se tiene un depósito $0$, clientes $1,\dots,n$ con demanda $q_i$, tiempo de servicio $s_i$ y ventana $[a_i,b_i]$, una flota de $K$ vehículos de capacidad $Q$ y tiempos de viaje $t_{ij}$. Se buscan rutas que empiecen y acaben en el depósito, visiten cada cliente exactamente una vez, no superen la capacidad y empiecen cada servicio dentro de su ventana (se puede esperar si se llega antes de $a_i$), minimizando el tiempo total de viaje.
\end{definicion}

Una formulación PLEM usa $x_{ij}\in\{0,1\}$ (un vehículo va de $i$ a $j$), el instante de inicio de servicio $T_i$ y la carga acumulada $u_i$:
\begin{align*}
\min\ &\sum_{i\ne j}t_{ij}x_{ij}\\
\st\ &\sum_{j}x_{ij}=1,\quad \sum_{j}x_{ji}=1\quad(i=1,\dots,n),\qquad \sum_j x_{0j}\le K,\\
&T_j\ge T_i+s_i+t_{ij}-M_{ij}(1-x_{ij})\quad(i\ne j,\ j\ge 1),\qquad a_i\le T_i\le b_i,\\
&u_j\ge u_i+q_j-Q(1-x_{ij})\quad(i\ne j,\ j\ge 1),\qquad q_i\le u_i\le Q .
\end{align*}
Las restricciones de tiempo y de carga eliminan además los subcircuitos, porque a lo largo de un ciclo sin depósito $T$ debería crecer indefinidamente.

\begin{proposicion}[Cota big-M de las ventanas]\label{prop:bigm-tw}
La restricción temporal con $M_{ij}=\max(0,\ b_i+s_i+t_{ij}-a_j)$ es válida: si $x_{ij}=1$ impone la llegada correcta y si $x_{ij}=0$ no elimina ninguna combinación de instantes que respete las ventanas.
\end{proposicion}

Con más de unas decenas de clientes, la formulación anterior es lenta, y los resolvedores de rutas usan búsqueda local sobre las propias rutas. OR-Tools modela el tiempo como una \emph{dimensión}: una variable acumulada por nodo, con una holgura que representa la espera y un rango que impone la ventana~\cite{ortools-vrptw}.

\addtocounter{definicion}{1}%

\section{Programación por restricciones y taller}

\begin{definicion}[Problema de satisfacción de restricciones]\label{def:csp}
Un \emph{problema de satisfacción de restricciones} consta de variables $v_1,\dots,v_k$ con dominios finitos $D_1,\dots,D_k$ y un conjunto de restricciones, cada una una relación sobre algunas variables. La \emph{programación por restricciones} (CP) lo resuelve alternando \emph{propagación} (eliminar de los dominios valores que no pueden formar parte de ninguna solución) y \emph{búsqueda} (fijar una variable y retroceder si se llega a un dominio vacío). Con un objetivo, se añade la restricción «objetivo mejor que el incumbente» cada vez que se encuentra una solución.
\end{definicion}

La CP brilla cuando las restricciones son lógicas o de secuenciación, difíciles de expresar linealmente. Sus restricciones globales encapsulan estructuras frecuentes con propagadores eficientes: \texttt{AllDifferent}, \texttt{NoOverlap} (intervalos que no se solapan en un recurso), \texttt{Cumulative} (recurso con capacidad). CP-SAT, de OR-Tools, combina propagación de restricciones con técnicas de resolvedores SAT y de programación entera, y exige coeficientes enteros~\cite{ortools-cpsat}.

\begin{definicion}[Job-shop]\label{def:jobshop}
En el problema de \emph{taller} (\emph{job-shop}) hay $J$ trabajos; el trabajo $j$ es una secuencia de operaciones $(m_{j1},p_{j1}),\dots,(m_{jk_j},p_{jk_j})$ que deben ejecutarse en ese orden, cada una en una máquina $m$ y con duración $p$. Cada máquina procesa una operación a la vez y sin interrupción. Se minimiza el \emph{makespan} $C_{\max}$, el instante en que termina el último trabajo.
\end{definicion}

En CP-SAT cada operación es un intervalo $[S_{jk},S_{jk}+p_{jk})$; se impone $S_{j,k+1}\ge S_{jk}+p_{jk}$, una restricción \texttt{NoOverlap} por máquina y $C_{\max}\ge S_{jk_j}+p_{jk_j}$~\cite{ortools-js}. En PLEM, cada par de operaciones en la misma máquina necesita una binaria de orden y dos restricciones big-M, y la relajación lineal resultante es muy débil.

\begin{proposicion}[Cotas inferiores del makespan]\label{prop:cota-ms}
Para todo programa factible,
\[
C_{\max}\ \ge\ \max\Bigl(\max_{m}\sum_{(j,k):\,m_{jk}=m}p_{jk},\ \ \max_j\sum_{k}p_{jk}\Bigr).
\]
\end{proposicion}

\addtocounter{definicion}{1}%

\begin{table}[h]
\centering\small
\caption{Tecnologías del tema~3: combinatoria, rutas y programación por restricciones.}\label{tab:tec3}
\begin{tabularx}{\textwidth}{@{}l>{\raggedright\arraybackslash}X>{\raggedright\arraybackslash}p{2.2cm}>{\raggedright\arraybackslash}X@{}}
\toprule
Herramienta & Para qué sirve & Licencia o tipo & Cuándo elegirla \\
\midrule
OR-Tools CP-SAT~\cite{ortools-cpsat} & CP con intervalos, \texttt{NoOverlap}, \texttt{Cumulative}; también PLEM con coeficientes enteros & abierta (Apache~2.0) & programación de talleres, turnos y calendarios \\
OR-Tools Routing~\cite{ortools-vrptw} & VRP con capacidades, ventanas, recogidas y entregas & abierta (Apache~2.0) & rutas reales de decenas a miles de paradas \\
PLEM (HiGHS, SCIP, Gurobi) & formulaciones con binarias de orden o de arco & abierta o comercial & instancias pequeñas, o con fuerte estructura lineal \\
Programación dinámica propia & exacto sobre subconjuntos & código propio & certificar instancias pequeñas \\
Heurísticas propias (2-opt, inserción) & solución rápida y controlable & código propio & tiempo real o reglas de negocio muy específicas \\
\bottomrule
\end{tabularx}
\end{table}

\addtocounter{definicion}{1}%


% =====================================================================
\chapter{Optimización continua y convexa}
\label{tema:4}
% =====================================================================

Cuando las variables son continuas y el problema es convexo, todo mínimo local es global, existen certificados de optimalidad y los multiplicadores tienen interpretación económica. Este tema sigue los capítulos~4 («Convex optimization problems») y~5 («Duality»), en especial las secciones~5.5 («Optimality conditions») y~5.6 («Perturbation and sensitivity analysis»), de Boyd y Vandenberghe~\cite{bv}, texto del curso EE364a de Stanford; el capítulo~12 de Nocedal y Wright~\cite{nw} para las condiciones de Karush--Kuhn--Tucker en problemas no convexos; y el capítulo~3 («Economic Dispatch of Thermal Units and Methods of Solution») de Wood, Wollenberg y Sheblé~\cite{wood} para el caso de estudio. Las herramientas son CVXPY~\cite{cvxpy,cvxpydoc} y SciPy. El módulo~7 (tema~1) introdujo la convexidad en el contexto del entrenamiento; aquí se añaden las restricciones.

\section{Problemas convexos}

\begin{definicion}[Problema convexo]\label{def:convexo}
El problema de la \cref{def:prob-opt} es \emph{convexo} si $X=\R^{n}$, $f$ y las $g_i$ son funciones convexas y las $h_j$ son afines, $h_j(x)=a_j\T x-b_j$. Su conjunto factible es entonces convexo, como intersección de subniveles de funciones convexas e hiperplanos.
\end{definicion}

\begin{teorema}[Local implica global]\label{teo:local-global}
En un problema convexo, todo óptimo local es óptimo global.
\end{teorema}

Las clases convexas más usadas, de menor a mayor generalidad, son la programación lineal, la cuadrática (objetivo $\tfrac12x\T Px+q\T x$ con $P\succeq 0$ y restricciones lineales), la cónica de segundo orden y la semidefinida. CVXPY sigue las reglas de la \emph{programación convexa disciplinada} (DCP): construye las expresiones con funciones atómicas de curvatura conocida y reglas de composición que garantizan la convexidad, de modo que rechaza un modelo cuya convexidad no puede certificar, y lo traduce a la forma cónica que aceptan resolvedores como Clarabel, SCS u OSQP~\cite{cvxpy,cvxpydoc}.

\section{Lagrangiana, dualidad y condiciones KKT}

\begin{definicion}[Lagrangiana y función dual]\label{def:lagrangiana}
Para el problema de la \cref{def:prob-opt} con $X=\R^{n}$, la \emph{lagrangiana} es
\[
\Lag(x,\lambda,\nu)=f(x)+\sum_{i=1}^{m}\lambda_ig_i(x)+\sum_{j=1}^{p}\nu_jh_j(x),\qquad\lambda\ge 0,
\]
y la \emph{función dual} es $q(\lambda,\nu)=\inf_x\Lag(x,\lambda,\nu)$.
\end{definicion}

\begin{proposicion}[Dualidad débil]\label{prop:dual-debil-conv}
Si $\lambda\ge 0$, entonces $q(\lambda,\nu)\le f(x)$ para todo $x$ factible; por tanto $q(\lambda,\nu)\le p^{*}$, el valor óptimo.
\end{proposicion}

En un problema convexo que cumple la \emph{condición de Slater} (existe un punto estrictamente factible, $g_i(x)<0$ para las restricciones no afines), la dualidad es fuerte: $\max q=p^{*}$ (Boyd y Vandenberghe~\cite{bv}, cap.~5).

\begin{teorema}[Condiciones de Karush--Kuhn--Tucker]\label{teo:kkt}
Sean $f$, $g_i$ y $h_j$ diferenciables. Se dice que $(x^{*},\lambda^{*},\nu^{*})$ cumple las condiciones KKT si
\begin{enumerate}
  \item \emph{factibilidad primal}: $g_i(x^{*})\le 0$, $h_j(x^{*})=0$;
  \item \emph{factibilidad dual}: $\lambda^{*}\ge 0$;
  \item \emph{holgura complementaria}: $\lambda_i^{*}g_i(x^{*})=0$ para todo $i$;
  \item \emph{estacionariedad}: $\nabla f(x^{*})+\sum_i\lambda_i^{*}\nabla g_i(x^{*})+\sum_j\nu_j^{*}\nabla h_j(x^{*})=0$.
\end{enumerate}
Si el problema es convexo y se cumplen las condiciones KKT, $x^{*}$ es óptimo global. Si el problema es convexo y cumple la condición de Slater, todo óptimo satisface las KKT para algún $(\lambda^{*},\nu^{*})$~\cite{bv}. En problemas no convexos, las KKT son necesarias en un mínimo local bajo una cualificación de restricciones, pero no suficientes~\cite{nw}.
\end{teorema}

\begin{proposicion}[Multiplicadores como sensibilidades]\label{prop:sensibilidad}
Sea el problema convexo perturbado $\min f(x)$ s.a. $g_i(x)\le u_i$, $h_j(x)=w_j$, con valor óptimo $p^{*}(u,w)$, y sea $(\lambda^{*},\nu^{*})$ dual óptimo del problema sin perturbar con dualidad fuerte. Entonces
\[
p^{*}(u,w)\ \ge\ p^{*}(0,0)-\lambda^{*\mathsf{T}}u-\nu^{*\mathsf{T}}w ,
\]
y si $p^{*}$ es diferenciable en $(0,0)$, $\partial p^{*}/\partial u_i=-\lambda_i^{*}$ y $\partial p^{*}/\partial w_j=-\nu_j^{*}$.
\end{proposicion}

Es la generalización de los precios sombra de la \cref{prop:sombra}: relajar una restricción activa abarata el óptimo, en primera aproximación, en su multiplicador por unidad.

\section{Caso de estudio: despacho económico}

\addtocounter{definicion}{1}%

\begin{observacion}[Cuando se pierde la convexidad]\label{obs:noconvexo}
El caso anterior deja de ser convexo con detalles realistas: costes con efecto de apertura de válvulas (términos sinusoidales), mínimos técnicos con arranque y parada (variables binarias, el problema de asignación de unidades del cap.~4 de~\cite{wood}) o mezclas con calidades que se multiplican por caudales (términos bilineales). Entonces las KKT son solo necesarias, los métodos locales (SQP, punto interior) pueden quedarse en óptimos locales y conviene arrancar desde varios puntos, usar un resolvedor global o reformular. En el ajuste de redes neuronales del módulo~7 el problema tampoco es convexo, pero allí basta con un buen mínimo local.
\end{observacion}

\begin{table}[h]
\centering\small
\caption{Tecnologías del tema~4: optimización continua.}\label{tab:tec4}
\begin{tabularx}{\textwidth}{@{}l>{\raggedright\arraybackslash}X>{\raggedright\arraybackslash}p{2.2cm}>{\raggedright\arraybackslash}X@{}}
\toprule
Herramienta & Para qué sirve & Licencia o tipo & Cuándo elegirla \\
\midrule
CVXPY~\cite{cvxpy} & modelización convexa con verificación DCP y duales & abierta (Apache~2.0) & cualquier problema convexo de tamaño medio \\
Clarabel, SCS, OSQP & resolvedores cónicos (Clarabel, SCS) y cuadráticos (OSQP) & abiertos (Apache~2.0; SCS, MIT) & los usa CVXPY; OSQP, para control predictivo rápido \\
SciPy \texttt{minimize} & SLSQP, \texttt{trust-constr}, L-BFGS-B & abierta (BSD-3) & no lineal local con pocas variables \\
Ipopt & punto interior no lineal a gran escala & abierta (EPL-2.0) & no lineal grande con derivadas (vía Pyomo o CasADi) \\
Gurobi & cuadrático y no convexo bilineal global & comercial & mezclas y redes con términos bilineales \\
\bottomrule
\end{tabularx}
\end{table}

\addtocounter{definicion}{1}%


% =====================================================================
\chapter{Metaheurísticas, optimización sin derivadas y multiobjetivo}
\label{tema:5}
% =====================================================================

Cuando el objetivo solo puede evaluarse (un simulador, un código heredado, un ensayo) o cuando la estructura del problema es tan irregular que ningún resolvedor exacto escala, se recurre a métodos que solo usan evaluaciones de $f$. Este tema sigue los capítulos~7 («Direct Methods»), 8 («Stochastic Methods»), 9 («Population Methods») y~12 («Multiobjective Optimization») de Kochenderfer y Wheeler~\cite{kw}, la clase~17 del curso 15.053~\cite{mit15053} para el recocido, y los artículos originales de Kirkpatrick, Gelatt y Vecchi~\cite{kirkpatrick} (recocido simulado), Hansen~\cite{hansen} (CMA-ES) y Deb \emph{et al.}~\cite{nsga2} (NSGA-II).

\section{Optimización de caja negra}

\begin{definicion}[Problema de caja negra]\label{def:caja-negra}
Un problema $\min_{x\in\mathcal{X}}f(x)$ es de \emph{caja negra} si solo se dispone de un oráculo que devuelve $f(x)$ (quizá con ruido) para cada $x$ consultado, sin gradientes ni estructura explícita. Su coste se mide por el \emph{número de evaluaciones}.
\end{definicion}

Hay dos regímenes. Si cada evaluación es barata (milisegundos), se pueden permitir miles o millones de evaluaciones y se usan metaheurísticas. Si es cara (horas de simulación o un experimento de planta), cada evaluación debe elegirse con cuidado con un modelo sustituto, como en la optimización bayesiana~\cite{garnett,frazier}, que la fase desarrolla más adelante. Los métodos directos locales, como Nelder--Mead o Powell (cap.~7 de~\cite{kw}), son adecuados cuando el problema es continuo, de pocas variables y con un único valle.

\section{Recocido simulado}

El recocido simulado~\cite{kirkpatrick} imita el enfriamiento lento de un metal. Desde una solución $x$ se propone un vecino $x'$ (por ejemplo, intercambiar dos departamentos) con un mecanismo simétrico, $\Prob(x\to x')=\Prob(x'\to x)$, y se acepta con la probabilidad de Metropolis
\[
\alpha(x,x')=\min\Bigl(1,\ e^{-(f(x')-f(x))/T}\Bigr),
\]
donde la temperatura $T>0$ disminuye según un calendario, típicamente $T_{k+1}=\rho T_k$ con $\rho\in(0{,}9;\ 0{,}999)$. Con $T$ alta se aceptan muchos empeoramientos y se explora; con $T$ baja solo se aceptan mejoras y se explota.

\begin{proposicion}[Distribución estacionaria de Boltzmann]\label{prop:boltzmann}
Sea $\mathcal{X}$ finito, la propuesta simétrica e irreducible y $T$ fija. La cadena de Markov del recocido satisface el equilibrio detallado con
\[
\pi_T(x)=\frac{e^{-f(x)/T}}{\sum_{z\in\mathcal{X}}e^{-f(z)/T}},
\]
que es por tanto una distribución estacionaria. Cuando $T\to 0$, $\pi_T$ se concentra uniformemente en el conjunto de mínimos globales.
\end{proposicion}

La distribución de Boltzmann es la del módulo~6 (tema~5). La garantía de convergencia al óptimo exige calendarios logarítmicos, demasiado lentos; en la práctica se usa el geométrico y varias ejecuciones independientes.

\section{Algoritmos genéticos y estrategias evolutivas}

Un \emph{algoritmo genético} mantiene una población de soluciones y la renueva con tres operadores (cap.~9 de~\cite{kw}): \emph{selección} de padres favoreciendo a los mejores (por ejemplo, por torneo entre dos individuos al azar), \emph{cruce} que combina dos padres y \emph{mutación} que perturba un hijo. Con \emph{elitismo} los mejores pasan intactos a la generación siguiente. La representación debe respetar las restricciones: para permutaciones, el \emph{cruce por orden} (OX) copia un tramo del primer padre y completa los huecos con los genes que faltan en el orden en que aparecen en el segundo, de modo que el hijo sigue siendo una permutación.

\begin{definicion}[CMA-ES]\label{def:cmaes}
La \emph{estrategia evolutiva con adaptación de la matriz de covarianza} (CMA-ES)~\cite{hansen} optimiza en $\R^{n}$ manteniendo una distribución $\mathcal N(m,\sigma^{2}C)$. En cada generación:
\begin{enumerate}[label=\arabic*.]
  \item muestrea $\lambda$ candidatos $x_k=m+\sigma y_k$, con $y_k\sim\mathcal N(0,C)$, y los ordena por $f$;
  \item actualiza la media con los $\mu$ mejores: $m\leftarrow\sum_{i=1}^{\mu}w_ix_{i:\lambda}$, con pesos positivos decrecientes que suman~1; la referencia propone $w_i\propto\ln\frac{\lambda+1}{2}-\ln i$;
  \item actualiza $C$ con los pasos que tuvieron éxito (actualizaciones de rango uno y de rango $\mu$) y ajusta $\sigma$ comparando la longitud del camino evolutivo con la esperada bajo selección aleatoria.
\end{enumerate}
El \emph{tamaño efectivo de selección} es $\mu_{\mathrm{eff}}=1/\sum_iw_i^{2}$.
\end{definicion}

CMA-ES aprende la forma del valle: en un problema mal condicionado, $C$ tiende a una matriz proporcional a la inversa del hessiano, de modo que se comporta como un método de segundo orden sin derivadas. Es invariante a transformaciones monótonas de $f$ (solo usa el orden) y a rotaciones del espacio. Es una elección sólida para calibrar simuladores de caja negra con decenas de parámetros continuos y miles de evaluaciones posibles.

\section{Optimización multiobjetivo}

\begin{definicion}[Dominancia y frente de Pareto]\label{def:pareto}
Sea $F=(f_1,\dots,f_k)$ un vector de objetivos que se minimizan. Se dice que $x$ \emph{domina} a $x'$, $x\succ x'$, si $f_i(x)\le f_i(x')$ para todo $i$ y $f_j(x)<f_j(x')$ para algún $j$. Un punto factible es \emph{Pareto-óptimo} si ningún otro punto factible lo domina. La imagen $F(\mathcal{X}^{\mathrm{P}})$ del conjunto de Pareto-óptimos es el \emph{frente de Pareto}.
\end{definicion}

\begin{proposicion}[Suma ponderada]\label{prop:suma-ponderada}
Si $w_i>0$ para todo $i$ y $x^{*}$ minimiza $\sum_iw_if_i(x)$ en $\mathcal{X}$, entonces $x^{*}$ es Pareto-óptimo. El recíproco es falso: los puntos del frente que no están en la frontera de su envolvente convexa no minimizan ninguna suma ponderada.
\end{proposicion}

El método de la $\varepsilon$-restricción, $\min f_1$ s.a. $f_2\le\varepsilon$, sí alcanza los puntos no apoyados variando $\varepsilon$ (\cref{eje:5-eps}). Para problemas de caja negra, el \emph{NSGA-II}~\cite{nsga2} combina un algoritmo genético con dos ideas:
\begin{enumerate}[label=(\alph*)]
  \item \emph{ordenación no dominada}: la población se divide en frentes $\mathcal{F}_1$ (no dominados), $\mathcal{F}_2$ (no dominados al quitar $\mathcal{F}_1$), etc., y el rango de un individuo es el índice de su frente;
  \item \emph{distancia de aglomeración} (\emph{crowding distance}): dentro de un frente, para cada objetivo se ordenan los puntos, los extremos reciben distancia infinita y cada punto interior suma $(f_k^{\text{sig}}-f_k^{\text{ant}})/(f_k^{\max}-f_k^{\min})$; se prefieren los puntos aislados para mantener la diversidad.
\end{enumerate}
La selección por torneo compara primero el rango y después la distancia, y la nueva generación se forma con los mejores frentes de la unión de padres e hijos (elitismo).

\addtocounter{definicion}{1}%

\addtocounter{figure}{1}%

\begin{table}[h]
\centering\small
\caption{Tecnologías del tema~5: metaheurísticas y multiobjetivo.}\label{tab:tec5}
\begin{tabularx}{\textwidth}{@{}l>{\raggedright\arraybackslash}X>{\raggedright\arraybackslash}p{2.2cm}>{\raggedright\arraybackslash}X@{}}
\toprule
Herramienta & Para qué sirve & Licencia o tipo & Cuándo elegirla \\
\midrule
\texttt{cma} (Hansen \emph{et al.}) & CMA-ES de referencia, con restricciones de caja y ruido & abierta (BSD-3) & calibrar simuladores continuos de caja negra \\
pymoo~\cite{pymoo} & NSGA-II, NSGA-III, indicadores de calidad de frentes & abierta (Apache~2.0) & problemas multiobjetivo con operadores estándar \\
SciPy \texttt{optimize} & recocido dual, evolución diferencial, Nelder--Mead, Powell & abierta (BSD-3) & continuo, pocas variables, sin dependencias extra \\
Optuna & búsqueda de hiperparámetros (muestreadores TPE y CMA-ES, multiobjetivo) & abierta (MIT) & ajustar modelos o configuraciones con evaluaciones caras \\
Código propio & recocido, búsqueda tabú, genéticos sobre una representación específica & propio & combinatoria con vecindarios propios del dominio \\
\bottomrule
\end{tabularx}
\end{table}

\addtocounter{definicion}{1}%


% =====================================================================
\chapter{Cuando la IA entra en la formulación: predecir y optimizar}
\label{tema:6}
% =====================================================================

En los temas anteriores los datos eran conocidos. En la práctica la demanda, los tiempos de viaje o los rendimientos se \emph{predicen}, y la decisión hereda el error de la predicción. Este tema estudia cómo combinar ambos pasos. Sigue el capítulo~1 («Introduction and Examples») de Birge y Louveaux~\cite{bl} y las transparencias de programación estocástica del curso EE364a de Stanford~\cite{ee364stoch}; el capítulo~17 («Optimization under Uncertainty») de Kochenderfer y Wheeler~\cite{kw}; los artículos de Elmachtoub y Grigas~\cite{spo} y de Bertsimas y Kallus~\cite{bk} sobre la integración de predicción y decisión; el de Bertsimas y Sim~\cite{bsim} y el libro de Ben-Tal, El Ghaoui y Nemirovski~\cite{btn} para la optimización robusta, y el de Smith y Winkler~\cite{smith} para la maldición del optimizador.

\section{Predecir y después optimizar}

\begin{definicion}[Predecir y después optimizar]\label{def:pto}
Sea un problema $\min_{x\in\mathcal{X}}\E[c(x,\xi)\mid w]$, donde $\xi$ es un parámetro incierto (la demanda) y $w$ el contexto observado (temperatura, festivos). El enfoque \emph{predecir y después optimizar} entrena un modelo $\hat\xi=m_\phi(w)$ y resuelve $\hat x(w)\in\argmin_{x\in\mathcal{X}}c(x,\hat\xi)$, sustituyendo el parámetro por su predicción puntual.
\end{definicion}

\begin{proposicion}[Cuándo basta la media]\label{prop:lineal}
Si el coste es lineal en el parámetro, $c(x,\xi)=\xi\T x$, y $\mathcal{X}$ no depende de $\xi$, la decisión que minimiza el coste esperado condicionado es la que resulta de sustituir $\xi$ por su media condicionada $\E[\xi\mid w]$.
\end{proposicion}

La proposición explica el éxito del enfoque cuando la incertidumbre está en los costes de un PL (rutas con tiempos inciertos y objetivo lineal) y su fracaso cuando el coste es no lineal en $\xi$: con inventario, el coste de quedarse corto y el de sobrar son distintos, y la media deja de ser suficiente. Incluso en el caso lineal, dos modelos con el mismo error cuadrático pueden inducir decisiones de calidad muy distinta: lo que importa no es el error en la predicción, sino el error en la decisión, que es lo que mide la pérdida SPO (\emph{smart predict-then-optimize}) de Elmachtoub y Grigas~\cite{spo} (\cref{eje:6-spo}). Bertsimas y Kallus~\cite{bk} proponen, en lugar de predecir un valor, aprender pesos sobre escenarios históricos parecidos al contexto y optimizar contra esa distribución empírica.

\section{El vendedor de periódicos}

\begin{definicion}[Problema del vendedor de periódicos]\label{def:newsvendor}
Se decide una cantidad $q\ge 0$ antes de conocer la demanda $D$, con función de distribución $F$. Cada unidad que falta cuesta $c_u>0$ (margen perdido) y cada unidad que sobra, $c_o>0$ (coste menos valor de saldo). El coste esperado es
\[
G(q)=\E\bigl[c_u(D-q)^{+}+c_o(q-D)^{+}\bigr].
\]
\end{definicion}

\begin{teorema}[Fractil crítico]\label{teo:fractil}
Si $F$ es continua, $q^{*}=F^{-1}\!\bigl(c_u/(c_u+c_o)\bigr)$ minimiza $G$. Si $D$ es discreta, es óptimo el menor $q$ con $F(q)\ge c_u/(c_u+c_o)$.
\end{teorema}

Si la demanda es normal, $q^{*}=\mu+\sigma\,\Phi^{-1}(c_u/(c_u+c_o))$: se produce por encima de la media cuando el margen perdido supera al coste del sobrante. Predecir la media, por precisa que sea, no basta; hace falta un \emph{cuantil}, y por tanto una estimación de la incertidumbre calibrada fuera de muestra.

\section{Programación estocástica y robusta}

\begin{definicion}[Programa estocástico en dos etapas]\label{def:dos-etapas}
En un \emph{programa en dos etapas} se fija una decisión $x$ antes de observar $\xi$ y, una vez observado, una decisión de recurso $y$:
\[
\mathrm{RP}=\min_{x\in\mathcal{X}}\ c\T x+\E_\xi\bigl[Q(x,\xi)\bigr],\qquad Q(x,\xi)=\min_{y\in\mathcal{Y}(x,\xi)}q(\xi)\T y .
\]
La \emph{aproximación por media muestral} (SAA) sustituye la esperanza por la media sobre $S$ escenarios $\xi^{1},\dots,\xi^{S}$ y, si todo es lineal, da un único PL con una copia de $y$ por escenario~\cite{ee364stoch}.
\end{definicion}

\begin{definicion}[Valor de la solución estocástica y de la información perfecta]\label{def:vss}
Para un problema de minimización, sea $\bar x$ la solución del \emph{problema del valor esperado} (sustituir $\xi$ por $\E\xi$), $\mathrm{EEV}=c\T\bar x+\E[Q(\bar x,\xi)]$ su coste esperado real, y $\mathrm{WS}=\E_\xi[\min_x(c\T x+Q(x,\xi))]$ el coste esperado si se conociera $\xi$ antes de decidir (\emph{wait-and-see}). Se definen el \emph{valor de la solución estocástica} $\mathrm{VSS}=\mathrm{EEV}-\mathrm{RP}$ y el \emph{valor esperado de la información perfecta} $\mathrm{EVPI}=\mathrm{RP}-\mathrm{WS}$~\cite{bl}.
\end{definicion}

\begin{proposicion}\label{prop:vss}
$\mathrm{WS}\le\mathrm{RP}\le\mathrm{EEV}$; en consecuencia, $\mathrm{VSS}\ge 0$ y $\mathrm{EVPI}\ge 0$.
\end{proposicion}

La \emph{optimización robusta} renuncia a la distribución y exige factibilidad o buen coste para todo $\xi$ de un \emph{conjunto de incertidumbre} $\Xi$: $\min_x\max_{\xi\in\Xi}c(x,\xi)$. Para restricciones lineales con coeficientes en una caja, el problema robusto sigue siendo lineal.

\begin{proposicion}[Contrapartida robusta de una caja]\label{prop:robusta}
Sea $x\ge 0$ y $a_j\in[\bar a_j-\delta_j,\ \bar a_j+\delta_j]$ con $\delta_j\ge 0$. La restricción $a\T x\le b$ se cumple para todos los $a$ de la caja si y solo si $(\bar a+\delta)\T x\le b$. Para $x$ de signo libre, la condición es $\bar a\T x+\delta\T\abs{x}\le b$, que se linealiza con la \cref{prop:abs}.
\end{proposicion}

La caja completa es muy conservadora: supone que todos los coeficientes toman su peor valor a la vez. Bertsimas y Sim~\cite{bsim} limitan a $\Gamma$ el número de coeficientes que se desvían y controlan así el \emph{precio de la robustez}: el coste que se paga por protegerse.

\begin{observacion}[La maldición del optimizador]\label{obs:maldicion}
Si se eligen la mejor de varias alternativas por su valor estimado, el valor estimado de la elegida sobrestima su valor real, aunque cada estimación sea insesgada: $\E[\max_i\hat J_i]\ge\max_i\E[\hat J_i]$~\cite{smith} (\cref{eje:6-maldicion}). Por eso el beneficio que promete un optimizador sobre datos predichos debe validarse con datos nuevos, nunca con los usados para elegir.
\end{observacion}

\section{Caso de estudio: plan de cerveza fresca con demanda predicha}

\addtocounter{definicion}{1}%

\begin{table}[h]
\centering\small
\caption{Tecnologías del tema~6: predecir y optimizar bajo incertidumbre.}\label{tab:tec6}
\begin{tabularx}{\textwidth}{@{}l>{\raggedright\arraybackslash}X>{\raggedright\arraybackslash}p{2.2cm}>{\raggedright\arraybackslash}X@{}}
\toprule
Herramienta & Para qué sirve & Licencia o tipo & Cuándo elegirla \\
\midrule
scikit-learn & regresión, cuantiles (\texttt{GradientBoostingRegressor} con pérdida cuantílica), validación temporal & abierta (BSD-3) & modelos de predicción tabulares \\
SciPy / HiGHS & PL de SAA con escenarios & abierta (BSD-3; MIT) & programas estocásticos lineales de tamaño medio \\
CVXPY~\cite{cvxpy} & modelos robustos y con restricciones probabilísticas convexas & abierta (Apache~2.0) & contrapartidas robustas cónicas, CVaR \\
Pyomo & modelos estocásticos por escenarios y descomposición & abierta (BSD-3) & problemas grandes en dos o más etapas \\
PyTorch & capas de optimización diferenciables para aprender con la pérdida de decisión & abierta (BSD) & enfoques SPO o de aprendizaje orientado a la decisión \\
\bottomrule
\end{tabularx}
\end{table}

\addtocounter{definicion}{1}%



\backmatter
% =====================================================================
\begin{thebibliography}{99}
\addcontentsline{toc}{chapter}{Bibliografía}
\markboth{Bibliografía}{}
\small

\bibitem{optimus}
A.~AhmadiTeshnizi, W.~Gao y M.~Udell.
\newblock OptiMUS: Scalable optimization modeling with (MI)LP solvers and large language models.
\newblock En \emph{Proceedings of the 41st International Conference on Machine Learning}, PMLR 235:577--596, 2024.
\newblock \url{https://proceedings.mlr.press/v235/ahmaditeshnizi24a.html}

\bibitem{btn}
A.~Ben-Tal, L.~El Ghaoui y A.~Nemirovski.
\newblock \emph{Robust Optimization}.
\newblock Princeton University Press, 2009.

\bibitem{bk}
D.~Bertsimas y N.~Kallus.
\newblock From predictive to prescriptive analytics.
\newblock \emph{Management Science}, 66(3), 2020. Prepublicación de 2014.
\newblock \url{https://arxiv.org/abs/1402.5481}

\bibitem{bsim}
D.~Bertsimas y M.~Sim.
\newblock The price of robustness.
\newblock \emph{Operations Research}, 52(1):35--53, 2004.

\bibitem{bt}
D.~Bertsimas y J.~N. Tsitsiklis.
\newblock \emph{Introduction to Linear Optimization}.
\newblock Athena Scientific, 1997. Texto del curso 6.251J del MIT.

\bibitem{bl}
J.~R. Birge y F.~Louveaux.
\newblock \emph{Introduction to Stochastic Programming}, 2.ª ed.
\newblock Springer, 2011. Cap.~1, «Introduction and Examples».

\bibitem{pymoo}
J.~Blank y K.~Deb.
\newblock pymoo: Multi-objective optimization in Python.
\newblock \emph{IEEE Access}, 8, 2020.
\newblock \url{https://arxiv.org/abs/2002.04504}

\bibitem{bv}
S.~Boyd y L.~Vandenberghe.
\newblock \emph{Convex Optimization}.
\newblock Cambridge University Press, 2004. Cap.~4, «Convex optimization problems», y cap.~5, «Duality». Texto del curso EE364a de Stanford.
\newblock \url{https://web.stanford.edu/~boyd/cvxbook/}

\bibitem{ee364stoch}
S.~Boyd y equipo docente de EE364a.
\newblock Stochastic programming. Transparencias complementarias del curso EE364a: Convex Optimization I.
\newblock Stanford University.
\newblock \url{https://web.stanford.edu/class/ee364a/lectures/stoch_prog.pdf}

\bibitem{amp}
S.~P. Bradley, A.~C. Hax y T.~L. Magnanti.
\newblock \emph{Applied Mathematical Programming}.
\newblock Addison-Wesley, 1977. Cap.~3, «Sensitivity Analysis», y cap.~9, «Integer Programming». Texto del curso 15.053 del MIT.
\newblock \url{https://web.mit.edu/15.053/www/}

\bibitem{nsga2}
K.~Deb, A.~Pratap, S.~Agarwal y T.~Meyarivan.
\newblock A fast and elitist multiobjective genetic algorithm: NSGA-II.
\newblock \emph{IEEE Transactions on Evolutionary Computation}, 6(2):182--197, 2002.

\bibitem{cvxpy}
S.~Diamond y S.~Boyd.
\newblock CVXPY: A Python-embedded modeling language for convex optimization.
\newblock \emph{Journal of Machine Learning Research}, 17(83):1--5, 2016.
\newblock \url{https://www.jmlr.org/papers/v17/15-408.html}

\bibitem{cvxpydoc}
Desarrolladores de CVXPY.
\newblock CVXPY User Guide: «Disciplined Convex Programming» y «Advanced Features» (variables duales).
\newblock \url{https://www.cvxpy.org/tutorial/index.html}

\bibitem{spo}
A.~N. Elmachtoub y P.~Grigas.
\newblock Smart «predict, then optimize».
\newblock \emph{Management Science}, 68(1):9--26, 2022.
\newblock \url{https://arxiv.org/abs/1710.08005}

\bibitem{deepmind}
R.~Evans y J.~Gao.
\newblock DeepMind AI reduces Google data centre cooling bill by 40\,\%.
\newblock Blog de Google DeepMind, 20 de julio de 2016.
\newblock \url{https://deepmind.google/discover/blog/deepmind-ai-reduces-google-data-centre-cooling-bill-by-40/}

\bibitem{frazier}
P.~I. Frazier.
\newblock A tutorial on Bayesian optimization.
\newblock Prepublicación, 2018.
\newblock \url{https://arxiv.org/abs/1807.02811}

\bibitem{garnett}
R.~Garnett.
\newblock \emph{Bayesian Optimization}.
\newblock Cambridge University Press, 2023. Edición libre en línea.
\newblock \url{https://bayesoptbook.com/}

\bibitem{ortools-cpsat}
Google OR-Tools.
\newblock CP-SAT Solver.
\newblock Documentación oficial.
\newblock \url{https://developers.google.com/optimization/cp/cp_solver}

\bibitem{ortools-js}
Google OR-Tools.
\newblock The Job Shop Problem.
\newblock Documentación oficial.
\newblock \url{https://developers.google.com/optimization/scheduling/job_shop}

\bibitem{ortools-vrptw}
Google OR-Tools.
\newblock Vehicle Routing Problem with Time Windows.
\newblock Documentación oficial.
\newblock \url{https://developers.google.com/optimization/routing/vrptw}

\bibitem{hansen}
N.~Hansen.
\newblock The CMA evolution strategy: A tutorial.
\newblock Prepublicación, 2016 (revisada en 2023).
\newblock \url{https://arxiv.org/abs/1604.00772}

\bibitem{highs}
HiGHS.
\newblock High performance software for linear optimization.
\newblock Documentación oficial; licencia MIT.
\newblock \url{https://highs.dev/}

\bibitem{orion}
C.~Holland, J.~Levis, R.~Nuggehalli, B.~Santilli y J.~Winters.
\newblock UPS optimizes delivery routes.
\newblock \emph{Interfaces}, 47(1):8--23, 2017.

\bibitem{hh}
Q.~Huangfu y J.~A.~J. Hall.
\newblock Parallelizing the dual revised simplex method.
\newblock \emph{Mathematical Programming Computation}, 10(1):119--142, 2018.

\bibitem{johnson}
S.~M. Johnson.
\newblock Optimal two- and three-stage production schedules with setup times included.
\newblock \emph{Naval Research Logistics Quarterly}, 1(1):61--68, 1954.

\bibitem{kirkpatrick}
S.~Kirkpatrick, C.~D. Gelatt y M.~P. Vecchi.
\newblock Optimization by simulated annealing.
\newblock \emph{Science}, 220(4598):671--680, 1983.

\bibitem{kw}
M.~J. Kochenderfer y T.~A. Wheeler.
\newblock \emph{Algorithms for Optimization}.
\newblock MIT Press, 2019. Caps.~1 «Introduction», 7 «Direct Methods», 8 «Stochastic Methods», 9 «Population Methods», 12 «Multiobjective Optimization» y~17 «Optimization under Uncertainty».

\bibitem{landdoig}
A.~H. Land y A.~G. Doig.
\newblock An automatic method of solving discrete programming problems.
\newblock \emph{Econometrica}, 28(3):497--520, 1960.

\bibitem{mit6251}
MIT OpenCourseWare.
\newblock 6.251J / 15.081J: Introduction to Mathematical Programming.
\newblock Massachusetts Institute of Technology, otoño de 2009. Clase~1 «Formulations», clases~9--12 (dualidad y sensibilidad) y~24--25 (optimización discreta), con lecturas de~\cite{bt}.
\newblock \url{https://ocw.mit.edu/courses/6-251j-introduction-to-mathematical-programming-fall-2009/}

\bibitem{mit15053}
J.~Orlin y E.~Nasrabadi.
\newblock 15.053: Optimization Methods in Management Science.
\newblock Massachusetts Institute of Technology, primavera de 2013. Clases~2 «Formulations of linear and non-linear programs», 6 «Sensitivity analysis and shadow prices», 10--11 (programación entera), 12 «Branch and bound», 13 «Cutting planes» y~17 «Traveling salesman problem».
\newblock \url{https://ocw.mit.edu/courses/15-053-optimization-methods-in-management-science-spring-2013/}

\bibitem{mit15053ps}
J.~Orlin y E.~Nasrabadi.
\newblock 15.053: Optimization Methods in Management Science. Hojas de problemas~4, 5 y~6.
\newblock Massachusetts Institute of Technology, primavera de 2013.
\newblock \url{https://ocw.mit.edu/courses/15-053-optimization-methods-in-management-science-spring-2013/pages/assignments/}

\bibitem{nw}
J.~Nocedal y S.~J. Wright.
\newblock \emph{Numerical Optimization}, 2.ª ed.
\newblock Springer, 2006. Cap.~12, teoría de la optimización con restricciones.

\bibitem{alphaevolve}
A.~Novikov, N.~V{\~u}, M.~Eisenberger, E.~Dupont \emph{et al.}
\newblock AlphaEvolve: A coding agent for scientific and algorithmic discovery.
\newblock Prepublicación, Google DeepMind, 2025.
\newblock \url{https://arxiv.org/abs/2506.13131}

\bibitem{pulp}
S.~Mitchell, A.~Kean, A.~Mason, M.~O'Sullivan, A.~Phillips y F.~Peschiera.
\newblock PuLP: documentación oficial (licencia MIT).
\newblock \url{https://coin-or.github.io/pulp/}

\bibitem{funsearch}
B.~Romera-Paredes, M.~Barekatain, A.~Novikov \emph{et al.}
\newblock Mathematical discoveries from program search with large language models.
\newblock \emph{Nature}, 625:468--475, 2024.

\bibitem{smith}
J.~E. Smith y R.~L. Winkler.
\newblock The optimizer's curse: Skepticism and postdecision surprise in decision analysis.
\newblock \emph{Management Science}, 52(3):311--322, 2006.

\bibitem{solomon}
M.~M. Solomon.
\newblock Algorithms for the vehicle routing and scheduling problems with time window constraints.
\newblock \emph{Operations Research}, 35(2):254--265, 1987.

\bibitem{sb}
R.~S. Sutton y A.~G. Barto.
\newblock \emph{Reinforcement Learning: An Introduction}, 2.ª ed.
\newblock MIT Press, 2018.
\newblock \url{http://incompleteideas.net/book/the-book-2nd.html}

\bibitem{wood}
A.~J. Wood, B.~F. Wollenberg y G.~B. Sheblé.
\newblock \emph{Power Generation, Operation, and Control}, 3.ª ed.
\newblock Wiley, 2014. Cap.~3, «Economic Dispatch of Thermal Units and Methods of Solution».

\bibitem{opro}
C.~Yang, X.~Wang, Y.~Lu, H.~Liu, Q.~V. Le, D.~Zhou y X.~Chen.
\newblock Large language models as optimizers.
\newblock Prepublicación, 2023.
\newblock \url{https://arxiv.org/abs/2309.03409}

\end{thebibliography}


\end{document}
